\documentclass[10pt,a4paper]{article}

% ------------------------------------------------
% PAQUETES
% ------------------------------------------------

\usepackage[utf8]{inputenc}
\usepackage[T1]{fontenc}
\usepackage[spanish]{babel}

\usepackage{graphicx}
\usepackage{geometry}
\usepackage{setspace}
\usepackage{amsmath}
\usepackage{amssymb}
\usepackage{url}
\usepackage{hyperref}
\usepackage{fancyhdr}
\usepackage{titlesec}
\usepackage{enumitem}
\usepackage{array}
\usepackage{booktabs}
\usepackage{float}

% ------------------------------------------------
% FORMATO DE TÍTULOS
% ------------------------------------------------

\titleformat{\section}
  {\normalfont\Large\bfseries\raggedright}
  {}
  {0pt}
  {}

\titleformat{\subsection}
  {\normalfont\large\bfseries\raggedright}
  {}
  {0pt}
  {}

% ------------------------------------------------
% CONFIGURACIÓN DE PÁGINA
% ------------------------------------------------

\geometry{
    left=2.5cm,
    right=2.5cm,
    top=2.0cm,
    bottom=2.2cm
}

% Interlineado
\setstretch{1.0}

% Espaciado entre párrafos
\setlength{\parindent}{0pt}
\setlength{\parskip}{0.25em}

% ------------------------------------------------
% ENCABEZADO
% ------------------------------------------------

\pagestyle{fancy}

\fancyhf{}

% Logo en la esquina superior izquierda
\fancyhead[L]{
    
\includegraphics[width=7.2cm]{informe/latex/logo-izquierda.png}
}

% Línea horizontal debajo del encabezado
\renewcommand{\headrulewidth}{0.5pt}

% Número de página
\fancyfoot[C]{\thepage}

% Espacio necesario para el encabezado
\setlength{\headheight}{1.25cm}

% ------------------------------------------------
% DOCUMENTO
% ------------------------------------------------

\begin{document}

% =================================================
% PORTADA
% =================================================

\begin{titlepage}

\thispagestyle{empty}

\begin{center}
\noindent

\includegraphics[width=10cm]{informe/latex/logo-izquierda.png}

\vspace{5.5cm}

\rule{\textwidth}{0.5pt}

\vspace{0.15cm}

{\Large\bfseries Propuesta SIG para una PYME\par}

\vspace{0.1cm}

{\normalsize ICN292 - Sistema de Información para la Gestión\par}

\vspace{0.15cm}

\rule{\textwidth}{0.5pt}

\vspace{1.2cm}

% ------------------------------------------------
% AUTORES
% ------------------------------------------------

\begin{center}
\renewcommand{\arraystretch}{1.25}
\setlength{\tabcolsep}{8pt}

\begin{tabular}{p{3.2cm}p{2.8cm}p{5.2cm}p{3.2cm}}
\hline
\textbf{Nombre} & \textbf{RUT} & \textbf{Rol en el equipo} & \textbf{Paralelo} \\ \hline

Felipe Salazar 
& 12345678-9 
& Modelador de datos (ER)
& 101 \\ \hline

Martín Barra 
& 12345678-9 
& Analista de procesos (BPMN)
& 101 \\ \hline

Montserratt Rojas 
& 21.497.666-8 
& Contacto con la PYME / Gestión
& 101 \\ \hline

Ignacia Álvarez 
& 20.675.835-K 
& Analista de requerimientos
& 101 \\ \hline

Trinidad Salinas 
& 21.435.964-2 
& Gestión del equipo y coordinación
& 101 \\ \hline

\end{tabular}
\end{center}

\vfill

\begin{center}

{\small\bfseries Repositorio GitHub\par}

\vspace{0.1cm}

{\small
\url{https://github.com/trinidad-usm/ICN292-Proyecto}
\par}

\vspace{0.5cm}

{\small Septiembre, 2026\par}

\vspace{0.5cm}

{\small Campus Vitacura\par}

\end{center}

\end{center}

\end{titlepage}
\raggedright

% =================================================
% INICIO DEL INFORME
% =================================================


% ------------------------------------------------
% ESTRUCTURA
% ------------------------------------------------

\section*{B. PYME y problema (evidencia)}

\subsection*{1. Identificación de la PYME}

\begin{itemize}[leftmargin=1.2cm]
    \item Nombre de fantasía: Alimentos Emy.
    \item Razón social: Comercializadora de Alimentos Eduardo Santibáñez Acuña EIRL.
    \item RUT de la empresa: 77.593.313-5.
    \item Rubro / giro comercial: venta al por mayor de huevos, lácteos, abarrotes y de otros alimentos N.C.P.; distribuidor oficial de productos lácteos de marcas como Quillayes/Surlat y Colun.
    \item Canal atendido: canal tradicional (minimarkets, almacenes de barrio y tiendas de necesidades básicas), bajo modalidad de "autoventa" (venta, facturación y despacho en el mismo momento de la visita al cliente).
    \item Zona geográfica / ubicación de bodega: Manuel Rodríguez 007, Lote B 01, Villa Olivar Andino, comuna de Olivar, Región de O'Higgins; atiende principalmente la zona de Rancagua y comunas cercanas (ej. Machalí, Doñihue).
    \item Tamaño: EIRL formalizada ante el SII, con 2 colaboradores fijos y 1 de apoyo; el propietario cumple simultáneamente el rol de vendedor, chofer, bodeguero, facturador y administrador.
    \item Antigüedad: aproximadamente 4 años de operación.
    \item Cartera de clientes: 45 clientes activos, atendidos mediante 3 rutas semanales fijas más rutas de fin de semana.
    \item Clasificación: empresa B2B ("última milla"), dedicada a la distribución mayorista hacia el canal tradicional.
\end{itemize}

\textbf{Datos del dueño / encargado (constancia y acreditación formal):}

\begin{itemize}[leftmargin=1.2cm]
    \item Nombre completo: Eduardo Santibáñez Acuña.
    \item RUT: 17.601.309-5.
    \item Cargo / rol formal: Fundador.
    \item Profesión / formación: Ingeniero Civil Industrial.
    \item Correo electrónico de contacto: congeladosemy@gmail.com.
    \item Teléfono de contacto: +56 9 8766 0844.
\end{itemize}


\subsection*{2. Evidencia de existencia}

Se sostuvo una reunión de levantamiento de información vía videollamada el 26 de agosto de 2026 con el propietario de la PYME, Eduardo (Alimentos Emy), en la que se explicó en detalle el modelo de negocio, los procesos actuales y las herramientas utilizadas (transcripción disponible como anexo).

\subsection*{3. Problema}

Alimentos Emy no cuenta con un proceso sistematizado de planificación de rutas de reparto ni de asignación del mix de carga por cliente y zona geográfica basado en su historial de ventas, lo que se traduce en un incumplimiento recurrente de la meta interanual de crecimiento en ventas (30\% respecto del mismo mes del año anterior) durante la temporada baja (marzo-agosto) y en sobrecostos de combustible, su principal partida de costo operativo.

\vspace{0.3cm}

\textbf{Evidencia cuantitativa disponible (declarada por el dueño):}
\begin{itemize}[leftmargin=1.2cm]
    \item Meta interna de crecimiento: 30\% de venta respecto al mismo mes del año anterior.
    \item Cumplimiento: meta alcanzada enero-mayo 2026 (con caídas en marzo), pero no alcanzada en junio y julio 2026; agosto 2026 proyectado a cerrar por debajo de la meta.  
    \item Merma: 0,7\% del volumen mensual de venta (declarada como controlada, no es el problema prioritario).
    \item Nivel de servicio (KPI logístico: \% de lo pedido por el cliente efectivamente despachado) no se mide hoy de forma sistemática en Alimentos Emy; el dueño reconoce que podría mejorar si tuviera mejor estimación de demanda por producto/zona.
    
\end{itemize}

\textbf{Impacto del problema:}
\begin{itemize}[leftmargin=1.2cm]
    \item Costo: el combustible es el principal costo operativo del negocio; sin una ruta y un mix de carga optimizados, se generan desplazamientos y viajes menos eficientes.
    \item Ventas: al no anticipar qué familias de producto demanda cada zona/cliente, se pierden oportunidades de venta (quiebres de stock en el camión) y se dificulta cumplir la meta de crecimiento. 
    \item Escalabilidad: toda la planificación de rutas y carga depende exclusivamente del criterio ("de guata", en palabras del dueño) de Eduardo, quien concentra todos los roles de la empresa; esto impide delegar la operación a un tercero o contratar personal de venta, limitando el crecimiento de la PYME.
    
\end{itemize}

\vspace{0.3cm}

En el formulario de levantamiento respondido directamente por el dueño, Eduardo describe el problema prioritario como el \textbf{aumento en la cartera de clientes mediante la optimización operativa}, y señala como impacto esperado la \textbf{reducción del costo logístico, el aumento en el número de atenciones diarias y el aumento de la rentabilidad del negocio}, lo que confirma y complementa el diagnóstico anterior.

\subsection*{4. Objetio del SIG propuesto}

Diseñar y modelar (a nivel de datos y proceso) un Sistema de Información para la Gestión que, a partir del historial de ventas ya digitalizado en el ERP de Alimentos Emy, apoye la planificación semanal de rutas y la asignación del mix de carga por cliente y zona geográfica, con el fin de: (i) optimizar el uso del vehículo y reducir el costo de combustible, (ii) aumentar el cumplimiento de la meta de crecimiento de ventas, y (iii) sentar las bases de información necesarias para delegar/tercerizar en el futuro parte de la operación de reparto.

Consultado directamente por este punto, el dueño confirma que la decisión/operación diaria que busca facilitar o automatizar principalmente con la propuesta es la \textbf{planificación logística}.

\newpage
\section*{C. Requerimientos desde cero}

\textbf{Actores y roles}

\begin{center}

\renewcommand{\arraystretch}{1.3}
\setlength{\tabcolsep}{6pt}

\begin{tabular}{|p{4cm}|p{10cm}|}
\hline

\textbf{Actor} & \textbf{Rol / relación con el sistema} \\ \hline

Dueño / vendedor / transportista (Eduardo) &
Usuario principal, confirmado por el propio dueño como "Fundador - Transportista". Revisa stock, arma el plan semanal, carga el camión, conduce, visita clientes, factura y despacha (autoventa). \\ \hline

Cliente (minimarket / almacén) &
Recibe la visita, realiza el pedido en el momento de la atención; no interactúa directamente con el sistema, pero es la fuente del dato de demanda. \\ \hline

ERP actual de la PYME (Bicom) &
Sistema externo, fuente de datos: historial de ventas, stock, facturación y clientes. El SIG propuesto se alimenta de sus reportes exportables (Excel). \\ \hline

WhatsApp (mensajería con clientes) &
Herramienta externa confirmada por el dueño para coordinar pedidos y avisos con los clientes; no está integrada al ERP ni a ningún sistema formal. \\ \hline

Vendedor/repartidor delegado (futuro) &
Rol proyectado a futuro: persona distinta del dueño que podría ejecutar una ruta siguiendo el plan generado por el sistema, sin depender del conocimiento tácito de Eduardo. \\ \hline

\end{tabular}

\end{center}

\textbf{Alcance (in / out)}

\vspace{0.3cm}

\textbf{Dentro del alcance (in):} confirmado por el dueño: planificación, ruteo y abastecimiento. En detalle:
\begin{itemize}[leftmargin=1.2cm]
    \item Ingesta y análisis del historial de ventas por cliente, producto/familia, zona y fecha, exportado desde el ERP actual.
    \item Agrupación de clientes por zona geográfica y frecuencia de visita.
    \item Generación de una sugerencia de ruta semanal y de mix de carga (cantidades por familia de producto) por ruta, respetando la capacidad del vehículo.
    \item Reporte/plan exportable para uso del vendedor (reemplazo del "planner" impreso actual).
    \item Seguimiento del cumplimiento de la meta de crecimiento de ventas (real vs. meta 30%).
\end{itemize}

\textbf{Fuera del alcance (out):} confirmado por el dueño: volumen de ventas, cobranza y gestiones comerciales. En detalle:
\begin{itemize}[leftmargin=1.2cm]
    \item Desarrollo de la aplicación de comparación de precios de retail para minimarkets (proyecto paralelo que la PYME está desarrollando con un tercero, mencionado en la entrevista como iniciativa distinta).
    \item Reemplazo del ERP/POS actual (facturación, boleteo, control de stock): el SIG propuesto se integra como apoyo analítico, no reemplaza esas funciones.
    \item Gestión de cobranza y de la relación comercial con clientes (negociación de precios, condiciones de pago).
    \item Optimización de rutas en tiempo real con tráfico/GPS (ruteo dinámico); se considera trabajo futuro.
\end{itemize}

Consultado por los requerimientos clave, el dueño confirma que la funcionalidad mínima que debe cumplir la herramienta para ser útil en terreno es un \textbf{plan de ruteo óptimo según demanda}, y que la principal condición de uso en terreno es la \textbf{facilidad de uso}.

\textbf{Requerimientos funcionales}

\begin{enumerate}[leftmargin=1.5cm, itemsep=0.4em]
    \item El sistema debe permitir importar el historial de ventas exportado desde el ERP (Excel/CSV) con al menos los campos cliente, producto, familia, cantidad y fecha. 
    \item El sistema debe calcular el volumen y la frecuencia de compra histórica de cada cliente, desagregado por familia de producto. 
    \item El sistema debe agrupar a los clientes por zona/sector geográfico a partir de su dirección o comuna.
    \item El sistema debe sugerir, para cada día de ruta, la lista de clientes a visitar según su frecuencia de compra asignada (semanal, quincenal, mensual). 
    \item El sistema debe sugerir el mix de carga (cantidad estimada por familia de producto) a llevar en el camión para una ruta determinada, en base al historial de compra de los clientes de esa ruta.
    \item El sistema debe permitir registrar la capacidad máxima de carga del vehículo (kg) como restricción del cálculo del mix de carga. 
    \item El sistema debe generar un reporte/plan de ruta y carga, exportable o imprimible, que reemplace el "planner" semanal actualmente confeccionado a mano.
    \item El sistema debe comparar la venta real mensual con la meta de crecimiento (30\% interanual) y mostrar el porcentaje de cumplimiento
    \item El sistema debe identificar clientes cuya frecuencia de compra reciente haya disminuido respecto de su patrón histórico, como alerta temprana de posible pérdida de cliente. 
    \item El sistema debe permitir, a futuro, que un vendedor distinto del dueño ejecute una ruta siguiendo el plan generado, sin requerir capacitación extensa sobre el criterio comercial de Eduardo. 
    
\end{enumerate}

\textbf{Requerimientos no funcionales}

\begin{enumerate}[leftmargin=1.5cm, itemsep=0.4em]
    \item Usabilidad: el sistema debe ser utilizable por una persona sin formación técnica (el propio dueño de la PYME).
    \item Seguridad y privacidad de datos: el sistema debe resguardar los datos comerciales y de clientes (nombres, direcciones, historial de compra) conforme a la Ley 21.719 sobre protección de datos personales, evitando exponerlos en repositorios públicos.
    \item Desempeño: la generación de la sugerencia de ruta y carga no debe tomar más de unos pocos minutos, de forma compatible con el hábito actual del dueño de planificar los domingos en la noche.
    \item Mantenibilidad: el sistema debe permitir actualizar periódicamente el historial de ventas (nueva exportación del ERP) sin requerir reconfigurar todo el modelo.
    \item Portabilidad de la salida: el plan de ruta y carga debe poder revisarse desde un dispositivo móvil o imprimirse, dado que el dueño hoy trabaja con un planner físico.
    \item Disponibilidad: el sistema debe estar disponible para consulta al menos los días domingo en la noche, momento en que el dueño arma su planificación semanal.

\end{enumerate}
\subsection*{Priorización de requerimientos (MoSCoW) y trazabilidad al problema}

A continuación se presenta la priorización de los requerimientos funcionales y no
funcionales bajo el esquema MoSCoW (\textit{Must have}, \textit{Should have},
\textit{Could have}, \textit{Won't have}), junto con su trazabilidad respecto del
problema de negocio y los objetivos del SIG definidos en la Sección B.

\subsubsection*{Requerimientos funcionales priorizados}

\begin{table}[H]
\centering
\small
\begin{tabular}{p{0.7cm} p{6.2cm} p{1.6cm} p{6.2cm}}
\toprule
\textbf{\#} & \textbf{Requisito} & \textbf{Prioridad} & \textbf{Trazabilidad (problema/objetivo)} \\
\midrule
RF1 & Importar historial de ventas (Excel/CSV: cliente, producto, familia, cantidad, fecha) & Must & Base de datos necesaria para toda la solución; sin ingesta no hay planificación posible. \\
RF2 & Calcular volumen y frecuencia de compra histórica por cliente/familia & Must & Objetivo (ii): aumentar el cumplimiento de la meta de crecimiento de ventas (30\,\%). \\
RF3 & Agrupar clientes por zona/sector geográfico & Must & Objetivo (i): optimizar el uso del vehículo y reducir el costo de combustible. \\
RF4 & Sugerir, por día de ruta, la lista de clientes a visitar según frecuencia asignada & Must & Objetivo (i) y (iii): estandarizar el criterio de ruta hoy dependiente del dueño. \\
RF5 & Sugerir el mix de carga (cantidad por familia de producto) por ruta & Must & Objetivo (i): reducir sobrecostos de combustible y quiebres de stock en el camión. \\
RF6 & Registrar la capacidad máxima de carga del vehículo (kg) como restricción & Must & Restricción física indispensable para que el ruteo sea factible (objetivo i). \\
RF7 & Generar reporte/plan de ruta y carga exportable o imprimible & Must & Reemplaza el \textit{planner} manual actual; habilita delegación futura (objetivo iii). \\
RF8 & Comparar venta real mensual con la meta de crecimiento (30\,\% interanual) & Should & Seguimiento directo del problema de incumplimiento de meta en temporada baja. \\
RF9 & Identificar clientes con caída en su frecuencia de compra histórica & Should & Alerta temprana de pérdida de clientes; apoya el objetivo de crecimiento en ventas. \\
RF10 & Permitir que un vendedor distinto del dueño ejecute una ruta siguiendo el plan generado & Could & Objetivo (iii): sentar bases para delegar/tercerizar la operación de reparto. \\
\bottomrule
\end{tabular}
\caption{Requerimientos funcionales priorizados según MoSCoW.}
\label{tab:rf-moscow}
\end{table}

\subsubsection*{Requerimientos no funcionales priorizados}

\begin{table}[H]
\centering
\small
\begin{tabular}{p{0.8cm} p{6.0cm} p{1.6cm} p{6.2cm}}
\toprule
\textbf{\#} & \textbf{Requisito} & \textbf{Prioridad} & \textbf{Trazabilidad (problema/objetivo)} \\
\midrule
RNF1 & Usabilidad: el sistema debe ser utilizable por una persona sin formación técnica & Must & El dueño concentra todos los roles de la empresa y sería el usuario principal del sistema. \\
RNF2 & Seguridad y privacidad de datos conforme a la Ley 21.719 & Must & Manejo de datos comerciales y de clientes (nombres, direcciones, historial de compra). \\
RNF3 & Desempeño: generación de la sugerencia de ruta y carga en pocos minutos & Should & Compatibilidad con el hábito actual del dueño de planificar los domingos en la noche. \\
RNF4 & Mantenibilidad: actualizar el historial de ventas sin reconfigurar todo el modelo & Should & Uso recurrente semanal del sistema con nuevas exportaciones del ERP. \\
RNF5 & Portabilidad de la salida: revisión desde dispositivo móvil o impresión & Should & El dueño hoy trabaja con un \textit{planner} físico impreso los domingos. \\
RNF6 & Disponibilidad para consulta al menos los domingos en la noche & Could & Momento específico en que el dueño arma la planificación semanal. \\
\bottomrule
\end{tabular}
\caption{Requerimientos no funcionales priorizados según MoSCoW.}
\label{tab:rnf-moscow}
\end{table}

\section*{D. Procesos BPMN}

\subsection*{1. BPMN As-Is del proceso crítico}

A continuación, se describe el flujo diario de trabajo de Alimentos Emy, desde la planificación matutina hasta el cierre de la jornada, con base en el formulario de levantamiento respondido directamente por el propietario, Eduardo, como complemento a la reunión del 26 de agosto de 2026. El proceso abarca la interacción entre el \textbf{Dueño / Operador}, el software de apoyo en terreno \textbf{ERP Bicom} y el \textbf{Cliente (Minimarket / Almacén)}.

\begin{figure}[htbp]
    \centering
    \includegraphics[width=\textwidth]{informe/latex/BPMN_as_is_pyme.png}
    \caption{Diagrama BPMN 2.0 del Proceso As-Is: Planificación, autoventa y cierre actual.}
    \label{fig:bpmn-as-is}
\end{figure}

\textbf{Descripción de las etapas del flujo actual:}

\begin{enumerate}[leftmargin=1.2cm]
    \item \textbf{Planificación y carga de camiones (inicio del día):} Según lo confirmado por el dueño, existe una demanda estimada mediante el conocimiento de los clientes, por lo que la decisión de qué productos y qué cantidades cargar al camión es por experiencia, sin datos concretos que la respalden. No se consulta ningún registro o reporte formal (por ejemplo, del ERP Bicom) antes de salir a ruta.
    \item \textbf{Secuencia y recorrido de la ruta:} El orden de visita a los clientes (Rancagua, Machalí, Doñihue) se define según los horarios de atención de cada cliente, y ante cualquier imprevisto en el camino la respuesta es mediante improvisación, sin un protocolo o criterio formal establecido.
    \item \textbf{Atención en el punto de venta (autoventa):} Los documentos de venta se generan en terreno mediante la aplicación Bicom; los datos de la venta se ingresan de manera semi automática al sistema en el momento de la atención, y el control de stock se realiza mediante inventarios físicos mensuales (no hay un chequeo de stock en tiempo real en cada visita).
    \item \textbf{Excepciones durante la ruta:} Si un cliente no está disponible, no tiene saldo para pagar o no requiere reposición ese día, la visita se reagenda o, en su defecto, se pierde de la ruta semanal (es decir, no siempre se recupera dentro del mismo ciclo).
    \item \textbf{Cierre de jornada y rendición (fin del día):} El cierre de caja se realiza mediante cierres diarios de efectivo, transferencias y pagos electrónicos, los cuales deben cuadrar con el total de ventas registradas en el sistema Bicom. Sin embargo, el dueño aclara explícitamente que \textbf{la mercadería no se cuadra de manera diaria}: el control físico del inventario del camión frente a lo facturado no ocurre día a día, sino que se apoya en los inventarios físicos mensuales mencionados en el punto 3, lo que constituye un punto de riesgo relevante para el control de merma y para el diseño del SIG.
\end{enumerate}

\subsection*{2. BPMN To-Be con la intervención del SIG}

El proceso propuesto integra un nuevo pool denominado \textbf{Sistema de Información para la Gestión (SIG)}, el cual centraliza el análisis de datos históricos para estandarizar la planificación logística y desacoplar la operación del criterio personal del fundador.

\begin{figure}[htbp]
    \centering
    \includegraphics[width=\textwidth]{informe/latex/BPMN_to_be_pyme.png}
    \caption{Diagrama BPMN 2.0 del Proceso To-Be: Intervención del SIG en planificación y control.}
    \label{fig:bpmn-to-be}
\end{figure}

\textbf{Secuencia del flujo To-Be intervenido:}
\begin{enumerate}[leftmargin=1.2cm]
    \item \textbf{Ingesta de datos y analítica predictiva:} El SIG importa el archivo histórico transaccional exportado desde el ERP Bicom. A partir de las frecuencias de compra y volumen por SKU, el sistema calcula la demanda proyectada para el ciclo semanal.
    \item \textbf{Generación asistida del plan logístico:} El SIG resuelve la optimización generando automáticamente: (i) la orden de carga sugerida por familia de productos sujeta al peso máximo del camión (kg), y (ii) la secuencia de ruta agrupada por cuadrante geográfico.
    \item \textbf{Carga y verificación en bodega:} El operador revisa la hoja de ruta y la lista de carga en el sistema (visualizable en formato digital/móvil) y realiza el carguío físico en bodega conforme al estándar analítico.
    \item \textbf{Autoventa guiada y registro de terreno:} El conductor efectúa el recorrido siguiendo la secuencia optimizada. La facturación in situ se mantiene en la app de Bicom, registrándose además las incidencias de visita para retroalimentar la base de clientes.
    \item \textbf{Cierre y actualización de métricas:} Al término del día se realiza la cuadratura financiera y el SIG actualiza de forma inmediata los indicadores clave (KPIs) de nivel de servicio, efectividad de ruta y cumplimiento de meta comercial.
\end{enumerate}

\subsection*{3. Breve explicación de mejoras}

La incorporación del SIG introduce mejoras directas en tres ejes operativos:

\begin{itemize}[leftmargin=1.2cm]
    \item \textbf{Qué se automatiza:}
    \begin{itemize}
        \item La consolidación del historial de facturación por cliente y familia de producto desde Bicom.
        \item El cálculo del mix de carga sugerido (kg y unidades) respetando la capacidad del furgón.
        \item La secuenciación ordenada del circuito de visitas semanal según zonas geográficas y periodicidad de compra.
    \end{itemize}
    \item \textbf{Qué se controla:}
    \begin{itemize}
        \item El límite físico de carga del vehículo para prevenir sobrepeso y viajes subutilizados.
        \item La periodicidad de visita a la cartera de 45 clientes, emitiendo alertas sobre cuentas con caídas en su patrón de reposición.
        \item La dependencia del conocimiento empírico del dueño, permitiendo que un tercero ejecute la autoventa siguiendo la hoja estandarizada.
    \end{itemize}
    \item \textbf{Qué se mide (KPIs):}
    \begin{itemize}
        \item \textbf{Nivel de servicio (Fill Rate):} Porcentaje de unidades demandadas por los almacenes que son efectivamente despachadas desde el camión sin incurrir en quiebre de stock[cite: 1].
        \item \textbf{Cumplimiento de la meta de venta:} Comparativa porcentual entre la venta real y la meta de crecimiento interanual del 30\%[cite: 1].
        \item \textbf{Eficiencia de ruta:} Clientes atendidos por jornada y reducción estimada de kilómetros recorridos por concepto de traslado.
    \end{itemize}
\end{itemize}

\section*{E. Modelo de datos preliminar}

\subsection*{1. Entidades centrales del proceso To-Be}

El modelo de datos preliminar se construyó a partir del proceso \textbf{BPMN To-Be}
(Figura 2), identificando las entidades de negocio que el Sistema de Información
para la Gestión (SIG) debe leer, transformar o generar en cada etapa del flujo
intervenido. Se priorizaron las entidades vinculadas directamente a los
\textit{data stores} del diagrama (``Historial y registro de ventas'') y a los
requerimientos funcionales definidos en la Sección C.

\begin{table}[H]
\centering
\small
\begin{tabular}{p{3.2cm} p{9.5cm}}
\toprule
\textbf{Entidad} & \textbf{Descripción y origen} \\
\midrule
\texttt{Cliente} & Minimarket/almacén atendido bajo modalidad de autoventa. Origen: RF3, RF4. \\
\texttt{Zona} & Sector geográfico de agrupación de clientes (Rancagua, Machalí, Doñihue). Origen: RF3. \\
\texttt{Producto} & SKU individual comercializado (ej. queso laminado 500\,g). Origen: RF1, RF5. \\
\texttt{FamiliaProducto} & Categoría de producto (quesos laminados, mantecosos, leche, etc.). Origen: RF2, RF5. \\
\texttt{Venta} & Encabezado de factura/boleta emitida en terreno vía Bicom ERP. Origen: data store ``Historial y registro de ventas''. \\
\texttt{DetalleVenta} & Línea de producto y cantidad asociada a una venta. Origen: RF1, RF2. \\
\texttt{Ruta} & Secuencia de visita semanal a un conjunto de clientes de una zona. Origen: RF4, RF7. \\
\texttt{Vehiculo} & Camión utilizado en la operación de reparto. Origen: RF6. \\
\texttt{PlanCarga} & Sugerencia de mix de carga (cantidad por familia) generada para una ruta. Origen: salida de la actividad ``Generación asistida del plan logístico''. \\
\texttt{AlertaCliente} & Registro de clientes con caída en su frecuencia de compra histórica. Origen: RF9. \\
\bottomrule
\end{tabular}
\caption{Entidades centrales del modelo de datos preliminar.}
\label{tab:entidades-er}
\end{table}

\subsection*{2. Atributos y clave primaria (PK) candidata}

\begin{table}[H]
\centering
\small
\begin{tabular}{p{2.8cm} p{7cm} p{2.9cm}}
\toprule
\textbf{Entidad} & \textbf{Atributos principales} & \textbf{PK candidata} \\
\midrule
\texttt{Cliente} & id\_cliente, nombre, dirección, id\_zona, frecuencia\_asignada & \texttt{id\_cliente} \\
\texttt{Zona} & id\_zona, nombre\_zona & \texttt{id\_zona} \\
\texttt{Producto} & sku, nombre, id\_familia & \texttt{sku} \\
\texttt{FamiliaProducto} & id\_familia, nombre\_familia & \texttt{id\_familia} \\
\texttt{Venta} & id\_venta, fecha, id\_cliente, monto\_total & \texttt{id\_venta} \\
\texttt{DetalleVenta} & id\_venta, sku, cantidad & \texttt{(id\_venta, sku)} \textit{compuesta} \\
\texttt{Ruta} & id\_ruta, día\_semana, id\_zona, id\_vehiculo & \texttt{id\_ruta} \\
\texttt{Vehiculo} & id\_vehiculo, capacidad\_kg & \texttt{id\_vehiculo} \\
\texttt{PlanCarga} & id\_ruta, id\_familia, cantidad\_sugerida & \texttt{(id\_ruta, id\_familia)} \textit{compuesta} \\
\texttt{AlertaCliente} & id\_cliente, fecha\_alerta, variación\_frecuencia & \texttt{(id\_cliente, fecha\_alerta)} \textit{compuesta} \\
\bottomrule
\end{tabular}
\caption{Atributos principales y clave primaria candidata por entidad.}
\label{tab:atributos-pk}
\end{table}

\subsection*{3. Relaciones y cardinalidades}

\begin{table}[H]
\centering
\small
\begin{tabular}{p{4cm} p{2cm} p{8cm}}
\toprule
\textbf{Relación} & \textbf{Cardinalidad} & \textbf{Justificación} \\
\midrule
Cliente -- Zona & N:1 & Cada cliente pertenece a una única zona geográfica de reparto. \\
Cliente -- Venta & 1:N & Un cliente acumula múltiples ventas a lo largo del tiempo (historial). \\
Venta -- DetalleVenta & 1:N & Cada venta (factura de autoventa) contiene una o más líneas de producto. \\
Producto -- FamiliaProducto & N:1 & Cada producto pertenece a una única familia (ej. queso laminado 500\,g $\rightarrow$ familia ``quesos laminados''). \\
Ruta -- Cliente & N:M & Una ruta visita a varios clientes, y un cliente puede aparecer en más de una ruta si su frecuencia lo requiere (quincenal, mensual). \\
Ruta -- Vehiculo & N:1 & Actualmente Eduardo opera con un solo vehículo para todas sus rutas semanales. \\
Ruta -- PlanCarga & 1:N & Cada ruta genera una sugerencia de carga compuesta por varias líneas, una por familia de producto. \\
Cliente -- AlertaCliente & 1:N & Un cliente puede generar más de una alerta en el tiempo si su frecuencia de compra sigue disminuyendo. \\
\bottomrule
\end{tabular}
\caption{Relaciones y cardinalidades del modelo ER preliminar.}
\label{tab:relaciones-er}
\end{table}

\subsection*{4. Trazabilidad proceso BPMN $\leftrightarrow$ datos}

La siguiente tabla justifica cómo el modelo entidad-relación sostiene directamente
al proceso BPMN To-Be (Figura 2), indicando qué entidad(es) se leen y cuál(es)
se escriben en cada actividad del flujo intervenido por el SIG.

\begin{table}[H]
\centering
\small
\begin{tabular}{p{4.3cm} p{4cm} p{4cm}}
\toprule
\textbf{Actividad BPMN To-Be} & \textbf{Entidad(es) que lee} & \textbf{Entidad(es) que escribe} \\
\midrule
Ingesta de datos y analítica predictiva & \texttt{Venta}, \texttt{DetalleVenta} & Demanda proyectada (insumo interno) \\
Generación asistida del plan logístico & \texttt{Cliente}, \texttt{Zona}, \texttt{Vehiculo} & \texttt{Ruta}, \texttt{PlanCarga} \\
Carga y verificación en bodega & \texttt{PlanCarga}, \texttt{Ruta} & --- (uso operativo, sin escritura) \\
Autoventa guiada y registro de terreno & \texttt{Ruta}, \texttt{Cliente} & \texttt{Venta}, \texttt{DetalleVenta} \\
Cierre y actualización de métricas & \texttt{Venta} vs. meta de crecimiento & \texttt{AlertaCliente}, indicadores KPI \\
\bottomrule
\end{tabular}
\caption{Trazabilidad entre el proceso BPMN To-Be y el modelo de datos.}
\label{tab:trazabilidad}
\end{table}

\subsection*{5. Justificación general}

El modelo entidad-relación propuesto sostiene el proceso BPMN To-Be en la medida
en que cada entidad corresponde directamente a un insumo o resultado de una
actividad del flujo intervenido. Las entidades \texttt{Venta} y
\texttt{DetalleVenta} materializan el \textit{data store} ``Historial y registro
de ventas'' presente de forma recurrente en el diagrama, y constituyen la base
sobre la cual el SIG calcula la demanda proyectada (RF1, RF2). Las entidades
\texttt{Cliente} y \texttt{Zona} permiten resolver la agrupación geográfica
requerida para la generación de rutas (RF3, RF4), mientras que
\texttt{Vehiculo} introduce la restricción de capacidad de carga (RF6) que
condiciona la generación de \texttt{PlanCarga} (RF5). Finalmente, la entidad
\texttt{AlertaCliente} traza directamente al requerimiento de detección temprana
de pérdida de clientes (RF9) y al problema de negocio descrito en la Sección B:
la falta de un proceso sistematizado de planificación que hoy depende
exclusivamente del criterio empírico (``de guata'') del dueño de la PYME.

% =================================================
% FIN DEL DOCUMENTO
% =================================================

\end{document}
